\documentclass[border=10pt]{standalone}
\usepackage[UTF8]{ctex}
\usepackage{tikz}
\usepackage{amsmath,amssymb}
\usetikzlibrary{arrows.meta,calc,positioning,fit,backgrounds}
\definecolor{cBlue}{RGB}{37,99,235}
\definecolor{cOrange}{RGB}{234,88,12}
\definecolor{cGray}{RGB}{100,116,139}
\definecolor{cGreen}{RGB}{22,163,74}
\definecolor{cRed}{RGB}{220,38,38}
\begin{document}
\begin{tikzpicture}[>=Stealth,line width=0.9pt,
   box/.style={draw,rounded corners=3pt,align=center,inner sep=6pt,line width=1pt,
               minimum width=3.5cm,minimum height=1.5cm},
   sbox/.style={draw,rounded corners=3pt,align=center,inner sep=4pt,line width=0.9pt,font=\small},
   lab/.style={align=center,font=\footnotesize}]

  % ---------- Row 1: the variable-compression pipeline ----------
  \node[box,draw=cBlue,fill=cBlue!4] (P) at (0,0)
     {一个 PDE\\[2pt]$\Delta(x,t,u,u_x,u_t,\dots)=0$\\[1pt]\footnotesize 多个自变量、解空间巨大};
  \node[box,draw=cOrange,fill=cOrange!6] (X) at (5.4,0)
     {对称生成元（旋钮）\\[2pt]$X=\xi\,\partial_x+\tau\,\partial_t+\phi\,\partial_u$\\[1pt]\footnotesize 作用在 $(x,t,u)$ 上的向量场 $=\mathfrak g$};
  \node[box,draw=cGreen,fill=cGreen!6] (O) at (10.8,0)
     {一个 ODE\\[2pt]在相似变量 $z$ 中\\[1pt]\footnotesize 少了一个自变量};
  \draw[->,cGray,line width=1.1pt] (P) -- (X) node[midway,above,lab,text=cGray]{求对称};
  \draw[->,cGreen,line width=1.2pt] (X) -- (O) node[midway,above,lab,text=cGreen]{沿对称约化};
  \node[lab,text=cGray] at (5.4,-1.45)
     {每拧一个对称“旋钮”，就消去一个自变量——这正是“把对称当变量压缩工具”};

  % ---------- Row 2: the Burgers concrete instance ----------
  \node[sbox,draw=cRed,fill=cRed!4] (B1) at (0,-3.5)
     {$u_t+u\,u_x=\nu\,u_{xx}$\\[1pt]\footnotesize 粘性 Burgers};
  \node[sbox,draw=cRed,fill=cRed!4,right=1.45cm of B1] (B2)
     {$z=x-ct$\\[1pt]$u=U(z)$};
  \node[sbox,draw=cRed,fill=cRed!4,right=1.2cm of B2] (B3)
     {$\nu U''=(U-c)\,U'$};
  \node[sbox,draw=cRed,fill=cRed!4,right=1.2cm of B3] (B4)
     {$U=c-k\tanh\dfrac{kz}{2\nu}$\\[1pt]\footnotesize 行波/激波解（精确）};
  \draw[->,cRed] (B1) -- (B2) node[midway,above,lab,text=cRed]{平移对称\\$\partial_t,\partial_x$};
  \draw[->,cRed] (B2) -- (B3) node[midway,above,lab,text=cRed]{代入};
  \draw[->,cRed] (B3) -- (B4) node[midway,above,lab,text=cRed]{解出};

  % ---------- Row 3: catalogue of canonical PDE-symmetry groups ----------
  \node[lab,text=cBlue] at (5.4,-5.1) {\textbf{典型 PDE 对称群（旋钮的种类）}};
  \node[lab,align=center] at (5.4,-6.05)
    {缩放 $(\mathbb{R}_{>0},\times)\ \to\ $相似解 / 量纲分析（Buckingham $\Pi$）；\quad
     伽利略 $\mathrm{Gal}\ \to\ $热、扩散方程\\[3pt]
     庞加莱 $ISO(1,3)\ \to\ $波动 / Klein--Gordon / Maxwell / Dirac；\quad
     共形 $\mathrm{Conf}\ \to\ $Laplace、波动\\[3pt]
     薛定谔 $\mathrm{Sch}=SL(2,\mathbb{R})\ltimes\mathrm{Gal}\ \to\ $自由热方程的最大对称};
\end{tikzpicture}
\end{document}
